% Options for packages loaded elsewhere
\PassOptionsToPackage{unicode}{hyperref}
\PassOptionsToPackage{hyphens}{url}
\PassOptionsToPackage{dvipsnames,svgnames,x11names}{xcolor}
\documentclass[
  10pt,
  fontset=fandol]{ctexart}
\usepackage{xcolor}
\usepackage[margin=16mm]{geometry}
\usepackage{amsmath,amssymb}
\setcounter{secnumdepth}{-\maxdimen} % remove section numbering
\usepackage{iftex}
\ifPDFTeX
  \usepackage[T1]{fontenc}
  \usepackage[utf8]{inputenc}
  \usepackage{textcomp} % provide euro and other symbols
\else % if luatex or xetex
  \usepackage{unicode-math} % this also loads fontspec
  \defaultfontfeatures{Scale=MatchLowercase}
  \defaultfontfeatures[\rmfamily]{Ligatures=TeX,Scale=1}
\fi
\usepackage{lmodern}
\ifPDFTeX\else
  % xetex/luatex font selection
\fi
% Use upquote if available, for straight quotes in verbatim environments
\IfFileExists{upquote.sty}{\usepackage{upquote}}{}
\IfFileExists{microtype.sty}{% use microtype if available
  \usepackage[]{microtype}
  \UseMicrotypeSet[protrusion]{basicmath} % disable protrusion for tt fonts
}{}
\makeatletter
\@ifundefined{KOMAClassName}{% if non-KOMA class
  \IfFileExists{parskip.sty}{%
    \usepackage{parskip}
  }{% else
    \setlength{\parindent}{0pt}
    \setlength{\parskip}{6pt plus 2pt minus 1pt}}
}{% if KOMA class
  \KOMAoptions{parskip=half}}
\makeatother
\usepackage{longtable,booktabs,array}
\newcounter{none} % for unnumbered tables
\usepackage{calc} % for calculating minipage widths
% Correct order of tables after \paragraph or \subparagraph
\usepackage{etoolbox}
\makeatletter
\patchcmd\longtable{\par}{\if@noskipsec\mbox{}\fi\par}{}{}
\makeatother
% Allow footnotes in longtable head/foot
\IfFileExists{footnotehyper.sty}{\usepackage{footnotehyper}}{\usepackage{footnote}}
\makesavenoteenv{longtable}
\setlength{\emergencystretch}{3em} % prevent overfull lines
\providecommand{\tightlist}{%
  \setlength{\itemsep}{0pt}\setlength{\parskip}{0pt}}
\usepackage{amsmath,amssymb,mathtools,needspace}
\xeCJKDeclareCharClass{CJK}{"2032,"0400->"04FF,"2160->"216B,"2460->"2473,"25A0,"25C7,"25CB,"25CF}
\IfFontExistsTF{Arial Unicode MS}{\xeCJKsetup{AutoFallBack=true}\setCJKfallbackfamilyfont{\CJKrmdefault}{Arial Unicode MS}}{}
\setcounter{secnumdepth}{0}
\setlength{\emergencystretch}{2em}
\allowdisplaybreaks
\usepackage{bookmark}
\IfFileExists{xurl.sty}{\usepackage{xurl}}{} % add URL line breaks if available
\urlstyle{same}
\hypersetup{
  pdftitle={Gauss 公式},
  colorlinks=true,
  linkcolor={Maroon},
  filecolor={Maroon},
  citecolor={Blue},
  urlcolor={Blue},
  pdfcreator={LaTeX via pandoc}}

\title{Gauss 公式}
\author{}
\date{}

\begin{document}
\maketitle

\subsection{4.4 Gauss 公式}\label{gauss-ux516cux5f0f}

形如式（4.1.3）的机械求积公式

（本句及公式续下页。）

\begin{quote}
校注（agent 补充，CH04B-E001）：原扫描式（4.3.17）首个分子明确印为 \(h^2\)，本转录原样保留；该处疑为原书排印错误。由本页式（4.3.16）及中点展开的梯形项相减，\(f''\) 项系数应为 \(h^3/12=h^3/(2!\times6)\)，才能在代入紧接着的 \(h\sum f''\) 关系后得到原页下一式的 \(h^2[f'(b)-f'(a)]/12\)。这不是 OCR 改写；\href{../assets/ch04b/review-p106-middle.png}{局部原图证据}。
\end{quote}

\href{../../source-images/pdf-107.jpeg}{原页图像}

\subsection{4.4 Gauss 公式（续）}\label{gauss-ux516cux5f0fux7eed}

\[
\int_a^b f(x)\,\mathrm{d}x\approx\sum_{k=0}^{n}A_kf(x_k)
\tag{4.4.1}
\]

中含有 \(2n+2\) 个待定参数 \(x_k,A_k\)（\(k=0,1,\cdots,n\)），适当选择这些参数，有可能使求积公式具有 \(2n+1\) 次代数精度。这类求积公式称为 Gauss 公式。

\subsubsection{4.4.1 Gauss 点}\label{gauss-ux70b9}

Gauss 公式的求积节点称为 Gauss 点。

\textbf{定义 4.3}　如果求积公式（4.4.1）具有 \(2n+1\) 次代数精度，则称其节点 \(x_k\)（\(k=0,1,\cdots,n\)）是 Gauss 点。

下面从分析 Gauss 点的特性着手研究 Gauss 公式的构造问题。

\textbf{定理 4.4}　对于插值型求积公式（4.4.1），其节点 \(x_k\)（\(k=0,1,\cdots,n\)）是 Gauss 点的充分必要条件，是以这些点为零点的多项式 \(\displaystyle\omega(x)=\prod_{k=0}^{n}(x-x_k)\) 与任意次数不超过 \(n\) 的多项式 \(P(x)\) 均正交，即

\[
\int_a^b P(x)\omega(x)\,\mathrm{d}x=0.
\tag{4.4.2}
\]

\textbf{证明}　先证必要性。设 \(P(x)\) 是任意次数不超过 \(n\) 的多项式，则 \(P(x)\omega(x)\) 的次数不超过 \(2n+1\)。因此，如果 \(x_0,x_1,\cdots,x_n\) 是 Gauss 点，则求积公式（4.4.1）对于 \(P(x)\omega(x)\) 能准确成立，即有

\[
\int_a^b P(x)\omega(x)\,\mathrm{d}x
=\sum_{k=0}^{n}A_kP(x_k)\omega(x_k).
\]

但 \(\omega(x_k)=0\)（\(k=0,1,\cdots,n\)），故式（4.4.2）成立。

再证充分性。对于任意给定的次数不超过 \(2n+1\) 的多项式 \(f(x)\)，用 \(\omega(x)\) 除 \(f(x)\)，记商为 \(P(x)\)，余式为 \(Q(x)\)，\(P(x)\) 与 \(Q(x)\) 都是次数不超过 \(n\) 的多项式：

\[
f(x)=P(x)\omega(x)+Q(x).
\]

而利用式（4.4.2），得

\[
\int_a^b f(x)\,\mathrm{d}x=\int_a^b Q(x)\,\mathrm{d}x.
\tag{4.4.3}
\]

由于所给求积公式（4.4.1）是插值型的，它对于 \(Q(x)\) 能准确成立，即

\[
\int_a^b Q(x)\,\mathrm{d}x=\sum_{k=0}^{n}A_kQ(x_k).
\]

再注意到 \(\omega(x_k)=0\)，知 \(Q(x_k)=f(x_k)\)，从而有

\[
\int_a^b Q(x)\,\mathrm{d}x=\sum_{k=0}^{n}A_kf(x_k),
\]

于是由式（4.4.3）知

\[
\int_a^b f(x)\,\mathrm{d}x=\sum_{k=0}^{n}A_kf(x_k).
\]

（证明续下页。）

\href{../../source-images/pdf-108.jpeg}{原页图像}

\subsubsection{4.4.1 Gauss 点（续）}\label{gauss-ux70b9ux7eed}

可见求积公式（4.4.1）对于一切次数不超过 \(2n+1\) 的多项式均能准确成立，因此 \(x_k\)（\(k=0,1,\cdots,n\)）是 Gauss 点。证毕。

\subsubsection{4.4.2 Gauss-Legendre 公式}\label{gauss-legendre-ux516cux5f0f}

不失一般性，可取 \(a=-1,b=1\) 而考察区间 \([-1,1]\) 上的 Gauss 公式①

\[
\int_{-1}^{1}f(x)\,\mathrm{d}x\approx\sum_{k=0}^{n}A_kf(x_k).
\tag{4.4.4}
\]

我们知道，Legendre 多项式（见式（3.4.1））是区间 \([-1,1]\) 上的正交多项式，因此，Legendre 多项式 \(P_{n+1}(x)\) 的零点就是求积公式（4.4.4）的 Gauss 点。形如式（4.4.4）的 Gauss 公式特别地称为 Gauss-Legendre 公式。

若取 \(P_1(x)=x\) 的零点 \(x_0=0\) 作节点构造求积公式

\[
\int_{-1}^{1}f(x)\,\mathrm{d}x\approx A_0f(0),
\]

令它对 \(f(x)=1\) 准确成立，即可定出 \(A_0=2\)。这样构造出的一点 Gauss-Legendre 公式是中矩形公式。

再取 \(P_2(x)=\dfrac12(3x^2-1)\) 的两个零点 \(\pm\dfrac1{\sqrt3}\) 构造求积公式：

\[
\int_{-1}^{1}f(x)\,\mathrm{d}x\approx
A_0f\left(-\frac1{\sqrt3}\right)+A_1f\left(\frac1{\sqrt3}\right),
\]

令它对于 \(f(x)=1,x\) 都准确成立，有

\[
\begin{cases}
A_0+A_1=2,\\[3pt]
A_0\left(-\dfrac1{\sqrt3}\right)+A_1\left(\dfrac1{\sqrt3}\right)=0.
\end{cases}
\]

由此解出 \(A_0=A_1=1\)，从而得到两点 Gauss-Legendre 公式

\[
\int_{-1}^{1}f(x)\,\mathrm{d}x\approx
f\left(-\frac1{\sqrt3}\right)+f\left(\frac1{\sqrt3}\right).
\]

三点 Gauss-Legendre 公式的形式是

\[
\begin{aligned}
&\int_{-1}^{1}f(x)\,\mathrm{d}x\\
&\quad\approx\frac59f\left(-\frac{\sqrt{15}}5\right)
+\frac89f(0)+\frac59f\left(\frac{\sqrt{15}}5\right).
\end{aligned}
\]

\textbf{表 4.5}

{\def\LTcaptype{none} % do not increment counter
\begin{longtable}[]{@{}rcc@{}}
\toprule\noalign{}
\(n\) & \(x_k\) & \(A_k\) \\
\midrule\noalign{}
\endhead
\bottomrule\noalign{}
\endlastfoot
0 & \(0.000\,000\,0\) & \(2.000\,000\,0\) \\
1 & \(\pm0.577\,350\,3\) & \(1.000\,000\,0\) \\
2 & \(\pm0.774\,596\,7\) & \(0.555\,555\,6\) \\
2 & \(0.000\,000\,0\) & \(0.888\,888\,9\) \\
3 & \(\pm0.861\,136\,3\) & \(0.347\,854\,8\) \\
3 & \(\pm0.339\,881\,0\) & \(0.652\,145\,2\) \\
4 & \(\pm0.906\,179\,3\) & \(0.236\,926\,9\) \\
4 & \(\pm0.538\,469\,3\) & \(0.478\,628\,7\) \\
4 & \(0.000\,000\,0\) & \(0.568\,888\,9\) \\
\end{longtable}
}

表 4.5 列出了 Gauss-Legendre 公式（4.4.4）的节点

（本句续下页。）

① 对于任意求积区间 \([a,b]\)，通过变换 \(x=\dfrac{b-a}{2}t+\dfrac{a+b}{2}\) 可以化到区间 \([-1,1]\) 上，这时

\[
\int_a^b f(x)\,\mathrm{d}x
=\frac{b-a}{2}\int_{-1}^{1}f\left(\frac{b-a}{2}t+\frac{a+b}{2}\right)\,\mathrm{d}t.
\]

\begin{quote}
校注（agent 补充，CH04B-E002）：表 4.5 的 \(n=3\) 第二对节点在原扫描中印为 \(\pm0.339\,881\,0\)，本表原样保留；这是疑似原书数值排印错误。\(P_4(x)=(35x^4-30x^2+3)/8\) 的较小正零点为 \(\sqrt{(15-2\sqrt{30})/35}\approx0.339\,981\,0436\)，保留七位小数应为 \(0.339\,981\,0\)。原表权重 \(0.652\,145\,2\) 未改动。\href{../assets/ch04b/review-table-4-5.png}{表 4.5 原扫描局部}。
\end{quote}

\href{../../source-images/pdf-109.jpeg}{原页图像}

\subsubsection{4.4.2 Gauss-Legendre 公式（续）}\label{gauss-legendre-ux516cux5f0fux7eed}

和系数。

\subsubsection{4.4.3 Gauss 公式的余项}\label{gauss-ux516cux5f0fux7684ux4f59ux9879}

\textbf{定理 4.5}　对于 Gauss 公式（4.4.1），其余项

\[
R(x)=\int_a^b f(x)\,\mathrm{d}x-\sum_{k=0}^{n}A_kf(x_k)
=\frac{f^{(2n+2)}(\xi)}{(2n+2)!}\int_a^b\omega^2(x)\,\mathrm{d}x,
\tag{4.4.5}
\]

这里 \(\omega(x)=(x-x_0)(x-x_1)\cdots(x-x_n)\)。

\textbf{证明}　以 \(x_0,x_1,\cdots,x_n\) 为节点构造次数不大于 \(2n+1\) 的多项式 \(H(x)\)，使其满足条件

\[
H(x_i)=f(x_i),\quad H'(x_i)=f'(x_i)\qquad(i=0,1,\cdots,n).
\]

我们知道，这样的 \(H(x)\) 称为 Hermite 插值多项式（见式（2.6.3）、（2.6.6））。由于 Gauss 公式具有 \(2n+1\) 次代数精度，它对于 \(H(x)\) 能准确成立，即

\[
\int_a^b H(x)\,\mathrm{d}x
=\sum_{k=0}^{n}A_kH(x_k)=\sum_{k=0}^{n}A_kf(x_k),
\]

故有

\[
\begin{aligned}
R(x)&=\int_a^b f(x)\,\mathrm{d}x-\sum_{k=0}^{n}A_kf(x_k)
=\int_a^b f(x)\,\mathrm{d}x-\int_a^b H(x)\,\mathrm{d}x\\
&=\int_a^b\frac{f^{(2n+2)}(\eta)}{(2n+2)!}\omega^2(x)\,\mathrm{d}x.
\end{aligned}
\]

再考虑到该函数 \(\omega^2(x)\) 在 \([a,b]\) 上保号，应用积分中值定理得式（4.4.5）。证毕。

\subsubsection{4.4.4 Gauss 公式的稳定性}\label{gauss-ux516cux5f0fux7684ux7a33ux5b9aux6027}

对比 Newton-Cotes 公式，Gauss 公式不但是高精度的，而且是数值稳定的。Gauss 公式的稳定性之所以能够得到保证，是由于它的求积系数具有非负性。

\textbf{定理 4.6}　Gauss 公式（4.4.1）的求积系数 \(A_k\)（\(k=0,1,\cdots,n\)）全是正的。

\textbf{证明}　考察

\[
l_k(x)=\prod_{\substack{j=0\\j\ne k}}^{n}\frac{x-x_j}{x_k-x_j},
\]

它们是 \(n\) 次多项式，因而 \(l_k^2(x)\) 是 \(2n\) 次多项式，故 Gauss 公式（4.4.1）对于它能准确成立，即有

\[
\int_a^b l_k^2(x)\,\mathrm{d}x=\sum_{i=0}^{n}A_i l_k^2(x_i).
\]

注意到 \(l_k(x_i)=\delta_{ki}\)，上式右端实际上即等于 \(A_k\)，从而有 \(\displaystyle A_k=\int_a^b l_k^2(x)\,\mathrm{d}x>0\)。证毕。

现在讨论求积过程的稳定性问题。在用求积公式 \(\displaystyle I_n=\sum_{k=0}^{n}A_kf(x_k)\) 进行实际计算时，通常不一定能提供准确的数据 \(f_k=f(x_k)\)，而只是给出含有误差（譬如由于计

（本句续下页。）

\href{../../source-images/pdf-110.jpeg}{原页图像}

\subsubsection{4.4.4 Gauss 公式的稳定性（续）}\label{gauss-ux516cux5f0fux7684ux7a33ux5b9aux6027ux7eed}

算机字长的限制引进的舍入误差）的数据 \(f_k^*\)，故实际求得的积分值为

\[
I_n^*=\sum_{k=0}^{n}A_kf_k^*.
\]

人们自然关心，数据误差对积分值的影响能否加以控制呢？

由于 Gauss 公式的求积系数具有非负性，故

\[
|I_n^*-I_n|\leq\sum_{k=0}^{n}A_k|f_k^*-f_k|
\leq\left(\sum_{k=0}^{n}A_k\right)\max_{0\leq k\leq n}|f_k^*-f_k|.
\]

再利用式（4.1.4）的第一式，有

\[
|I_n^*-I_n|\leq(b-a)\max_{0\leq k\leq n}|f_k^*-f_k|,
\]

由此即可断定 Gauss 公式是稳定的。

\subsubsection{4.4.5 带权的 Gauss 公式}\label{ux5e26ux6743ux7684-gauss-ux516cux5f0f}

考察积分 \(\displaystyle I=\int_a^b\rho(x)f(x)\,\mathrm{d}x\)，这里 \(\rho(x)\geq0\) 称为权函数，当 \(\rho(x)\equiv1\) 时即为普通的积分。

可以仿照处理普通积分的方法讨论带权的积分。考察求积公式

\[
\int_a^b\rho(x)f(x)\,\mathrm{d}x\approx\sum_{k=0}^{n}A_kf(x_k),
\]

如果它对于任意次数不超过 \(2n+1\) 的多项式均能准确地成立，则称之为 Gauss 型的。上述 Gauss 公式的求积节点 \(x_k\) 仍称为 Gauss 点。同样地，\(x_k\)（\(k=0,1,\cdots,n\)）是 Gauss 点的充要条件为，\(\displaystyle\omega(x)=\prod_{k=0}^{n}(x-x_k)\) 是区间 \([a,b]\) 上关于权函数 \(\rho(x)\) 的正交多项式。

若 \(a=-1,b=1\)，且取权函数 \(\rho(x)=\dfrac1{\sqrt{1-x^2}}\)，则所建立的 Gauss 公式为

\[
\int_{-1}^{1}\frac{f(x)}{\sqrt{1-x^2}}\,\mathrm{d}x
\approx\sum_{k=0}^{n}A_kf(x_k).
\tag{4.4.6}
\]

式（4.4.6）称为 Gauss-Chebyshev 公式。由于区间 \([-1,1]\) 上关于权函数 \(\dfrac1{\sqrt{1-x^2}}\) 的正交多项式是 Chebyshev 多项式（见 3.4.2 节），因此求积公式（4.4.6）的 Gauss 点是 \(n+1\) 次 Chebyshev 多项式的零点，即为

\[
x_k=\cos\left(\frac{2k+1}{2n+2}\pi\right)\qquad(k=0,1,\cdots,n).
\]

值得指出的是，运用正交多项式的零点构造 Gauss 求积公式，这种方法只是针对某些特殊的权函数才有效。构造 Gauss 公式的一般方法则是 4.1.2 节介绍过的待定系数法。

例如，设要构造下列形式的 Gauss 公式：

（公式续下页。）

\href{../../source-images/pdf-111.jpeg}{原页图像}

\subsubsection{4.4.5 带权的 Gauss 公式（续）}\label{ux5e26ux6743ux7684-gauss-ux516cux5f0fux7eed}

\[
\int_0^1\sqrt{x}f(x)\,\mathrm{d}x\approx A_0f(x_0)+A_1f(x_1),
\tag{4.4.7}
\]

令它对于 \(f(x)=1,x,x^2,x^3\) 准确成立，得

\[
\begin{cases}
A_0+A_1=\dfrac23,\\[5pt]
x_0A_0+x_1A_1=\dfrac25,\\[5pt]
x_0^2A_0+x_1^2A_1=\dfrac27,\\[5pt]
x_0^3A_0+x_1^3A_1=\dfrac29.
\end{cases}
\tag{4.4.8}
\]

由于

\[
x_0A_0+x_1A_1=x_0(A_0+A_1)+(x_1-x_0)A_1,
\]

利用式（4.4.8）的第一式，可将第二式化为

\[
\frac23x_0+(x_1-x_0)A_1=\frac25.
\]

同样地，利用第二式化第三式，利用第三式化第四式，分别得

\[
\frac25x_0+(x_1-x_0)x_1A_1=\frac27,\qquad
\frac27x_0+(x_1-x_0)x_1^2A_1=\frac29.
\]

从上面三个式子中消去 \((x_1-x_0)A_1\)，有

\[
\begin{cases}
\dfrac25x_0+\left(\dfrac25-\dfrac23x_0\right)x_1=\dfrac27,\\[5pt]
\dfrac27x_0+\left(\dfrac27-\dfrac25x_0\right)x_1^2=\dfrac29,
\end{cases}
\]

即

\[
\begin{cases}
\dfrac25(x_0+x_1)-\dfrac23x_0x_1=\dfrac27,\\[5pt]
\dfrac27(x_0+x_1)-\dfrac25x_0x_1=\dfrac29.
\end{cases}
\]

由此解出

\[
x_0x_1=\frac5{21},\quad x_0+x_1=\frac{10}9,
\]

从而求出

\[
\begin{aligned}
x_0&=0.821\,162,\quad x_1=0.289\,949,\\
A_0&=0.389\,111,\quad A_1=0.277\,556.
\end{aligned}
\]

于是形如式（4.4.7）的 Gauss 公式是

\[
\int_0^1\sqrt{x}f(x)\,\mathrm{d}x
\approx0.389\,111f(0.821\,162)+0.277\,556f(0.289\,949).
\]

\begin{quote}
校注（agent 补充，CH04B-E003）：原扫描在``从上面三个式子中消去''后的第二行印为 \(\dfrac27x_0+\left(\dfrac27-\dfrac25x_0\right)x_1^2=\dfrac29\)，本转录原样保留。疑误是括号后 \(x_1^2\) 的指数：前一行有 \((x_1-x_0)x_1A_1=\dfrac27-\dfrac25x_0\)，代入再上一组的第四式应只再乘 \(x_1\)，故此处应为 \(\left(\dfrac27-\dfrac25x_0\right)x_1\)，才与紧接着``即''后的第二行相符。这是原书疑似排印错误，非 OCR 错误。\href{../assets/ch04b/review-p111-elimination.png}{局部原扫描}。
\end{quote}

\subsection{本节覆盖原页的校注记录}\label{ux672cux8282ux8986ux76d6ux539fux9875ux7684ux6821ux6ce8ux8bb0ux5f55}

下列记录按原页关联，含同页相邻小节的校注；计算时同时读取上文校注和完整原页。

\begin{itemize}
\tightlist
\item
  \texttt{CH04B-E001}：原书 PDF 页 106
\item
  \texttt{CH04B-E002}：原书 PDF 页 108
\item
  \texttt{CH04B-E003}：原书 PDF 页 111
\end{itemize}

\end{document}
